\begin{figure}[H]
    \centering
    \resizebox{\textwidth}{!}{
    \begin{tikzpicture}[
        table/.style={
            rectangle split,
            rectangle split parts=2,
            draw=black!70,
            thick,
            font=\small,
            align=left,
            rounded corners=3pt,
            rectangle split part fill={gray!15, white}
        },
        fato/.style={
            table,
            font=\small\bfseries
        },
        dimensao/.style={
            table,
            font=\footnotesize
        },
        relation/.style={
            -latex,
            very thick,
            draw=black!50
        },
        background_box/.style={
            rectangle, draw=gray!40, thick, dashed, fill=gray!3, rounded corners=6pt, inner sep=15pt
        },
        label_text/.style={font=\scriptsize\sffamily\bfseries, fill=gray!5, draw=gray!40, rounded corners=2pt, inner sep=3pt}
    ]

    % --- FATOS (EM CIMA) ---
    \node[fato] (fato_focos) {
        FATO\_FOCOS\_INCENDIO
        \nodepart{two}
        \normalfont
        PK: foco\_id \\
        \textit{timestamp} \\
        M: satelite \\
        M: latitude, longitude \\
        M: frp (Potência Radiativa) \\
        M: risco\_ons\_calc
    };

    \node[fato, right=3cm of fato_focos] (fato_eventos) {
        FATO\_EVENTOS\_INCENDIO
        \nodepart{two}
        \normalfont
        PK: evento\_id \\
        \textit{timestamp\_inicio} \\
        \textit{timestamp\_fim} \\
        M: densidade \\
        M: risco\_cascata\_global \\
        M: area\_estimada\_km2
    };

    % --- DIMENSOES (EM BAIXO) ---
    % Linha 1 de Dimensões (4 tabelas)
    \node[dimensao, below left=3cm and -1.5cm of fato_focos] (dim_mun) {
        \textbf{DIM\_MUNICIPIOS}
        \nodepart{two}
        PK: codigo\_ibge \\
        estado\_sigla \\
        municipios\_nome
    };

    \node[dimensao, right=0.6cm of dim_mun] (dim_ti) {
        \textbf{DIM\_TERRAS\_INDIGENAS}
        \nodepart{two}
        PK: codigo\_ti \\
        nome\_terra\_indigena \\
        area\_hectares
    };

    \node[dimensao, right=0.6cm of dim_ti] (dim_hosp) {
        \textbf{DIM\_HOSPITAIS}
        \nodepart{two}
        PK: codigo\_cnes \\
        nome\_hospital \\
        tipo\_unidade
    };

    \node[dimensao, right=0.6cm of dim_hosp] (dim_rod) {
        \textbf{DIM\_RODOVIAS}
        \nodepart{two}
        PK: id\_trecho\_rodoviario \\
        codigo\_rodovia \\
        jurisdicao
    };

    % Linha 2 de Dimensões (3 tabelas centralizadas)
    \node[dimensao, below=1.5cm of dim_ti, xshift=-1cm] (dim_lt) {
        \textbf{DIM\_LINHAS\_TRANSMISSAO}
        \nodepart{two}
        PK: id\_linha\_transmissao \\
        nome\_linha \\
        tensao\_kv
    };

    \node[dimensao, right=0.6cm of dim_lt] (dim_sub) {
        \textbf{DIM\_SUBESTACOES}
        \nodepart{two}
        PK: id\_subestacao \\
        nome\_subestacao \\
        capacidade\_mw
    };
    
    \node[dimensao, right=0.6cm of dim_sub] (dim_sis) {
        \textbf{DIM\_SISTEMAS\_ISOLADOS}
        \nodepart{two}
        PK: id\_sistema\_isolado \\
        nome\_sistema\_isolado \\
        prioridade
    };

    % --- GRUPO BACKGROUND (DIMENSÕES COMPARTILHADAS) ---
    \begin{scope}[on background layer]
        \coordinate (bg_top) at ([yshift=0.8cm]dim_mun.north);
        \node[background_box, fit=(bg_top) (dim_mun) (dim_ti) (dim_hosp) (dim_rod) (dim_lt) (dim_sub) (dim_sis)] (bg_dimensoes) {};
        \node[anchor=north west, font=\bfseries\color{black!70}, xshift=5pt, yshift=-5pt] at (bg_dimensoes.north west) {Dimensões Compartilhadas (Esquema Constelação)};
    \end{scope}

    % --- RELACIONAMENTOS (Setas Reduzidas) ---
    % Apenas 2 setas direcionadas ao Container de Dimensões
    \draw[relation] (fato_focos.south) -- node[label_text, right=3pt] {Foreign Keys} (bg_dimensoes.north -| fato_focos.south);
    \draw[relation] (fato_eventos.south) -- node[label_text, left=3pt] {Foreign Keys} (bg_dimensoes.north -| fato_eventos.south);

    \end{tikzpicture}
    }
    \caption{Modelo Entidade-Relacionamento (Constelação).}
    \fonte{Elaborada pelos autores.}
    \label{fig:diagrama_er}
\end{figure}
